En este informe se analizan distintos algoritmos de ordenamiento (insertion sort, merge sort, quick sort, sort y panda sort) y de multiplicación de matrices (naive y Strassen), evaluando su desempeño en tiempo de ejecución y uso de memoria. Los experimentos se realizaron en un entorno controlado con datasets de diferentes tamaños y tipos de entrada. Los resultados muestran que los algoritmos eficientes mantienen un rendimiento estable frente a distintas condiciones, mientras que los de complejidad cuadrática y Strassen presentan mayores costos en tiempo o memoria en los casos estudiados. Estos hallazgos resaltan la importancia de contrastar la teoría con la práctica al seleccionar algoritmos para problemas concretos.